\documentclass[11pt,a4paper]{article}

\usepackage[utf8]{inputenc}

\usepackage[T1]{fontenc}

\usepackage[polish]{babel}

\usepackage{lmodern}

\usepackage{geometry}

\geometry{margin=27mm}

\usepackage[hidelinks]{hyperref}

\title{Polityki decyzyjne zależne od odpowiedzialności sprawcy}

\author{Mateusz Pomianek}

\date{Szkic naukowy, wersja 1.0.0 | 28 listopada 2026}

\begin{document}

\maketitle

\begin{abstract}

Proponowana miara główna to różnica prawdopodobieństw wyboru sankcji w H1.

\end{abstract}

\section{Wprowadzenie}

W ocenie moralnej informacja o intencji i skutku może mieć różną wagę zależnie od ocenianej reakcji (Cushman, 2008). Sankcję, wsparcie pokrzywdzonego i wsparcie sprawcy należy zatem rejestrować oddzielnie.

\section{Cel i pytania badawcze}

Hipotezy robocze wymagają zatwierdzenia przed rozpoczęciem właściwej weryfikacji.

\paragraph{H1.} Przy tej samej wielkości szkody działanie zamierzone zwiększy prawdopodobieństwo wyboru sankcji w porównaniu z działaniem nieumyślnym.

\paragraph{H2.} Możliwość rzeczywistego zadośćuczynienia zwiększy wybór rozwiązania naprawczego, a efekt będzie większy w wariancie nieumyślnym.

\section{Materiał i metody}

Pierwszorzędowy wskaźnik: różnica prawdopodobieństw wyboru sankcji w H1.

Losowy przydział kombinacji intencji, szkody i możliwości naprawy pozwoli oszacować kontrasty w modelu mieszanym. Główny kontrast dotyczy decyzji o sankcji w H1. Współczucie dla obu stron należy mierzyć niezależnie od decyzji.

Przed badaniem z udziałem ludzi należy zatwierdzić wymagania etyczne, zasady zgody i ochronę danych.

\section{Plan analizy}

Zamrożenie estymandu, reguł oceny jakości i obsługi braków nastąpi przed dostępem do danych wynikowych. Raport obejmie oszacowanie kontrastu głównego z przedziałem niepewności i pełne mianowniki. Jeśli przewidziano hipotezy drugorzędowe, przed analizą należy ustalić sposób kontroli wielokrotności.

\section{Wyniki}

W niniejszym pliku nie umieszczono wyników empirycznych, statystyk ani nieprzeprowadzonych analiz. Sekcja wymaga uzupełnienia po uzyskaniu i weryfikacji materiału dowodowego.

\section{Dyskusja i ograniczenia}

Ograniczenia obejmują możliwe błędy pomiaru, dobór próby, alternatywne mechanizmy i ograniczoną możliwość uogólniania. Korelacja nie dowodzi przyczynowości; poprawność techniczna nie jest dowodem znaczenia klinicznego.

\section{Wnioski}

Hipotezy i sposób pomiaru stanowią podstawę oceny wykonalności oraz przygotowania właściwego protokołu.

\section*{Bibliografia}

{\setlength{\parindent}{0pt}
\setlength{\parskip}{5pt}
\hangindent=1.27cm
Cushman, F. (2008). Crime and punishment: Distinguishing the roles of causal and intentional analyses in moral judgment. \textit{Cognition, 108}(2), 353--380. \url{https://doi.org/10.1016/j.cognition.2008.03.006}\par
}

\end{document}
